\documentclass[12pt,a4paper,reqno]{amsart}

% language
\usepackage[greek,english]{babel}
\usepackage[utf8]{inputenc}
\usepackage{alphabeta}

% change default names to greek
\addto\captionsenglish{
    \renewcommand{\contentsname}{Περιεχόμενα}
    \renewcommand{\refname}{Βιβλιογραφία}
    \renewcommand{\datename}{Ημερομηνία:}
    \renewcommand{\urladdrname}{Ιστοσελίδα}
}

% math 
\usepackage{amsmath,amsthm,amssymb,amscd}

% font
\usepackage[cal=euler]{mathalfa}
\usepackage{libertinus-type1}
% \usepackage{txfonts} % for upright greek letters
\usepackage{bm} % for bold symbols
\usepackage{bbm} % for the simply-looking bb symbols

% miscellaneous 
\usepackage{changepage} %for indenting environments
\usepackage{csquotes} % example: \textcquote{}
\usepackage{blkarray}
\setcounter{MaxMatrixCols}{20} % default for pmatrix is 10!!
\usepackage{ytableau}
\usepackage{array} %needed to increase the vertical length in a tabular

% drawing
\usepackage{tikz,tikz-cd}
\usetikzlibrary{shapes.misc, patterns, matrix, calc, intersections,positioning,arrows.meta}
\usepackage{graphics,graphicx}
\usepackage{float} % provides enhanced control and customization options for floating objects such as figures and tables

% colors
\usepackage{xcolor}
\definecolor{darkcandyapplered}{rgb}{0.64, 0.0, 0.0}
\definecolor{midnightblue}{rgb}{0.1, 0.1, 0.44}
\definecolor{mylightblue}{HTML}{336699}
\definecolor{burntorange}{rgb}{0.8, 0.33, 0.0}
\definecolor{iceberg}{rgb}{0.44, 0.65, 0.82}
\definecolor{applegreen}{rgb}{0.55, 0.71, 0.0}
\definecolor{canaryyellow}{rgb}{1.0, 0.94, 0.0}
\definecolor{lightgray}{rgb}{0.83, 0.83, 0.83}

% hrefs
\usepackage{hyperref}
\usepackage[noabbrev,capitalize]{cleveref}
\hypersetup{
    pdftoolbar=true,        
    pdfmenubar=true,        
    pdffitwindow=false,     
    pdfstartview={FitH},  % fits the width of the page to the window
    pdftitle={},
    pdfauthor={},
    pdfsubject={},
    pdfkeywords={},
    pdfnewwindow=true,  % links in new window
    colorlinks=true,  % false: boxed links; true: colored links
    linkcolor=darkcandyapplered,   % color of internal links
    citecolor=midnightblue,  % color of links to bibliography
    urlcolor=cyan,  % color of external links
    linktocpage=true  % changes the links from the section body to the page number
    }

% geometry
\textwidth=16cm 
\textheight=21cm 
\hoffset=-55pt 
\footskip=25pt

% thm envs (you might need to change the path)
\theoremstyle{definition}
\newtheorem{exercise}{Άσκηση}
\newtheorem{solution}{Λύση}
\newtheorem*{remark}{Παρατήρηση}
\newtheorem*{example}{Παράδειγμα}

% math ops 
% fields
\newcommand{\NN}{\mathbbmss{N}} 
\newcommand{\ZZ}{\mathbbmss{Z}} 
\newcommand{\QQ}{\mathbbmss{Q}} 
\newcommand{\RR}{\mathbbmss{R}} 
\newcommand{\CC}{\mathbbmss{C}} 
\newcommand{\KK}{\mathbbmss{K}} 
\newcommand{\FF}{\mathbbmss{F}} 
\newcommand{\HH}{\mathbbmss{H}} 

% symmetric group
\newcommand{\fS}{\mathfrak{S}}  
\newcommand{\fp}{\mathfrak{p}}  
\newcommand{\fm}{\mathfrak{m}}  

% calligraphic 
\newcommand{\aA}{\mathcal{A}} 
\newcommand{\bB}{\mathcal{B}}
\newcommand{\cC}{\mathcal{C}}
\newcommand{\dD}{\mathcal{D}}
\newcommand{\eE}{\mathcal{E}}
\newcommand{\fF}{\mathcal{F}}
\newcommand{\hH}{\mathcal{H}}
\newcommand{\iI}{\mathcal{I}}
\newcommand{\lL}{\mathcal{L}}
\newcommand{\nN}{\mathcal{N}}
\newcommand{\oO}{\mathcal{O}}
\newcommand{\pP}{\mathcal{P}}
\newcommand{\sS}{\mathcal{S}}
\newcommand{\mM}{\mathcal{M}}
\newcommand{\uU}{\mathcal{U}}

% bold
\newcommand{\bfa}{\mathbf{a}}
\newcommand{\bfe}{\mathbf{e}}
\newcommand{\bfF}{\pmb{F}}
\newcommand{\bfR}{\pmb{R}}
\newcommand{\bfv}{\mathbf{v}}
%\newcommand{\bfx}{\bm{x}}
%\newcommand{\bfx}{\mathbf{x}} 
\newcommand{\bfx}{\pmb{x}}
\newcommand{\bfX}{\pmb{X}}
\newcommand{\bfy}{\pmb{y}}
\newcommand{\bfz}{\pmb{z}}

% roman
\newcommand{\rmA}{\mathrm{A}}
\newcommand{\rmB}{\mathrm{B}}
\newcommand{\rmC}{\mathrm{C}}
\newcommand{\rmD}{\mathrm{D}} 
\newcommand{\rmH}{\mathrm{H}} 
\newcommand{\rmh}{\mathrm{h}} 
\newcommand{\rmI}{\mathrm{I}} 
\newcommand{\rmJ}{\mathrm{J}} 
\newcommand{\rmK}{\mathrm{K}}
\newcommand{\rmM}{\mathrm{M}}
\newcommand{\rmP}{\mathrm{P}}  
\newcommand{\rmp}{\mathrm{p}}  
\newcommand{\rmQ}{\mathrm{Q}}  
\newcommand{\rmR}{\mathrm{R}}
\newcommand{\rmS}{\mathrm{S}}
\newcommand{\rmT}{\mathrm{T}}
\newcommand{\rmU}{\mathrm{U}}
\newcommand{\rmV}{\mathrm{V}}
\newcommand{\rmY}{\mathrm{Y}}
\newcommand{\rmZ}{\mathrm{Z}}
\newcommand{\rmz}{\mathrm{z}}

% arrows and symbols 
\renewcommand{\to}{\rightarrow}
\newcommand{\toto}{\longrightarrow}
\newcommand{\mapstoto}{\longmapsto}
\newcommand{\then}{\Rightarrow}
\newcommand{\IFF}{\Leftrightarrow}
\newcommand{\tl}{\tilde}
\newcommand{\wtl}{\widetilde}
\newcommand{\ol}{\overline}
\newcommand{\ul}{\underline}
\newcommand{\oldemptyset}{\emptyset}
\renewcommand{\emptyset}{\varnothing}
\DeclareMathSymbol{\Arg}{\mathbin}{AMSa}{"39} % for arguments 
\newcommand{\onto}{\ensuremath{\twoheadrightarrow}}
\newcommand{\tle}{\trianglelefteq}
\newcommand{\tge}{\trianglerighteq}

% custom ops
\newcommand{\rmKer}{\mathrm{Ker}}
\newcommand{\rmIm}{\mathrm{Im}}
\newcommand{\lcm}{\mathrm{lcm}}
\newcommand{\GL}{\mathrm{GL}}
\newcommand{\ev}{\mathrm{ev}}
\newcommand{\End}{\mathrm{End}}
\newcommand{\rmid}{\mathrm{id}}
\newcommand{\rmspan}{\mathrm{span}}
\newcommand{\Spec}{\mathrm{Spec}}
\newcommand{\mSpec}{\mathrm{mSpec}}
\newcommand{\rad}{\mathrm{rad}}
\newcommand{\Rad}{\mathrm{Rad}}
\newcommand{\Jac}{\mathrm{Jac}}
\newcommand{\tor}{\mathrm{tor}}
\newcommand{\Ann}{\mathrm{Ann}}
\newcommand{\Hom}{\mathrm{Hom}}
\newcommand{\Hilb}{\mathrm{Hilb}}

% absolute value symbol
\usepackage{mathtools} 
\DeclarePairedDelimiter\abs{\lvert}{\rvert}%
\DeclarePairedDelimiter\norm{\lVert}{\rVert}%
\makeatletter
\let\oldabs\abs
\def\abs{\@ifstar{\oldabs}{\oldabs*}}

% tensor symbol
\newcommand{\tensor}[1]{%
  \mathbin{\mathop{\otimes}\limits_{#1}}%
}

% permutation cycle notation
\ExplSyntaxOn
\NewDocumentCommand{\cycle}{ O{\;} m }
 {
  (
  \alec_cycle:nn { #1 } { #2 }
  )
 }

\seq_new:N \l_alec_cycle_seq
\cs_new_protected:Npn \alec_cycle:nn #1 #2
 {
  \seq_set_split:Nnn \l_alec_cycle_seq { , } { #2 }
  \seq_use:Nn \l_alec_cycle_seq { #1 }
 }
\ExplSyntaxOff

% setminus symbol
\newcommand{\mysetminusD}{\hbox{\tikz{\draw[line width=0.6pt,line cap=round] (3pt,0) -- (0,6pt);}}}
\newcommand{\mysetminusT}{\mysetminusD}
\newcommand{\mysetminusS}{\hbox{\tikz{\draw[line width=0.45pt,line cap=round] (2pt,0) -- (0,4pt);}}}
\newcommand{\mysetminusSS}{\hbox{\tikz{\draw[line width=0.4pt,line cap=round] (1.5pt,0) -- (0,3pt);}}}
\newcommand{\sm}{\mathbin{\mathchoice{\mysetminusD}{\mysetminusT}{\mysetminusS}{\mysetminusSS}}}

% miscellaneous commands
\newcommand{\defn}[1]{{\color{mylightblue}{#1}}}
\newcommand{\toDo}{{\bf\color{red} TODO}}
\newcommand{\toCite}{{\bf\color{green} CITE}}
\newcommand*{\vertbar}{\rule[-1ex]{0.5pt}{2.5ex}} % for matrices with column vectors
\newcommand*{\horzbar}{\rule[.5ex]{2.5ex}{0.5pt}} % for matrices with row vectors
\newcommand{\myblue}[1]{{\color{iceberg}{#1}}}
\newcommand{\myorange}[1]{{\color{burntorange}{#1}}}
\newcommand{\mygreen}[1]{{\color{applegreen}{#1}}}
\newcommand{\myred}[1]{{\color{darkcandyapplered}{#1}}}

% ferrer's diagram
\newcommand{\fdiagram}[1]{
    \begin{tikzpicture}[scale=.7]
        \fill foreach \Z [count=\Y] in {#1}
        {foreach \X in {1,...,\Z} 
        {(\X,-\Y) circle[radius=3pt]}};
    \end{tikzpicture}
}

%
\usepackage{pgfplots}
\pgfplotsset{compat=1.18}

%
\makeatletter
\def\mathcolor#1#{\@mathcolor{#1}}
\def\@mathcolor#1#2#3{%
  \protect\leavevmode
  \begingroup
    \color#1{#2}#3%
  \endgroup
}
\makeatother

% 
\newenvironment{nouppercase}{%
  \let\uppercase\relax%
  \renewcommand{\uppercasenonmath}[1]{}}{}

% titlepage
\title{226: Θεωρία Δακτυλίων και Modules \\ Ασκήσεις με Ενδεικτικές Λύσεις}
\author[Β.~Δ. Μουστακας]{Βασίλης Διονύσης Μουστάκας \\ Πανεπιστήμιο Κρήτης}
\date{22 Απριλίου 2026}
% \urladdr{\href{https://sites.google.com/view/vasmous}{https://sites.google.com/view/vasmous}}

\begin{document}

\begingroup
\def\uppercasenonmath#1{} % this disables uppercase title
\let\MakeUppercase\relax % this disables uppercase authors
\maketitle
\endgroup

\setcounter{exercise}{28}
% \setcounter{theorem}{5}
\setcounter{equation}{7}
\renewcommand{\thesubsection}{\thesection.\arabic{subsection}}

Σε ότι ακολουθεί υποθέτουμε ότι $R$ είναι ένας μεταθετικός δακτύλιος, $F$ είναι ένα σώμα και $n$ είναι ένας θετικός ακέραιος, εκτός αν αναφέρεται διαφορετικά.

\begin{exercise}{\rm(\cite[Άσκηση~7.2.3]{DF04})}
  \label{ex:DF_7.2.3}
  Θεωρούμε το σύνολο $R^{\NN}$ όλων των (άπειρων) ακολουθιών της μορφής $(a_0, a_1, \dots)$, όπου $a_0, a_1, \dots \in R$ και ορίζουμε 
  \begin{align*}
    (a_n)_{n \in \NN} + (b_n)_{n \in \NN} &= (a_n + b_n)_{n \in \NN} \\
    (a_n)_{n \in \NN}\cdot(b_n)_{n \in \NN} &= (c_n)_{n \in \NN},
  \end{align*}
  πρόσθεση και πολλαπλασιασμό, αντίστοιχα, όπου 
  \[
  c_n = a_0b_n + a_1b_{n-1} + \cdots + a_nb_0 = \sum_{k + \ell = n} a_kb_\ell.
  \]
  \begin{itemize}
    \item[(α)] Να δείξετε ότι ο $\RR^\NN$ με μονάδα $(1, 0, 0, \dots)$ και μηδέν $(0, 0, \dots)$ είναι μεταθετικός δακτύλιος.
  \end{itemize}
  Συμβολίζουμε ένα αυθαίρετο στοιχείο αυτού του δακτύλιου με 
  \[
  \sum_{n \ge 0} a_n x^n = a_0 + a_1x + \cdots 
  \]
  και γράφουμε $R[\![x]\!]$ αντί για $R^\NN$. Το στοιχείο αυτό ονομάζεται \defn{τυπική δυναμοσειρά} (formal power series) και ο $R[\![x]\!]$ ονομάζεται δακτύλιος των τυπικών δυναμοσειρών. Να δείξετε τα εξής.
  \begin{itemize}
    \item[(β)] Η τυπική δυναμοσειρά 
    \[
    \sum_{n\ge0}x^n = 1 + x + x^2 + \cdots 
    \] 
    είναι αντιστρέψιμο στοιχείο του $R[\![x]\!]$ και να βρείτε το αντίστροφό της.
    \item[(γ)] Έχουμε $\sum_{n\ge0} a_nx^n \in \left(
      R[\![x]\!]
    \right)^\times$ αν και μόνο αν $a_0 \in R^\times$. 
  \end{itemize}
\end{exercise}

\begin{proof}[Λύση]
  Το (α) αφήνεται ως άσκηση. Για το (β) παρατηρούμε ότι 
  \[
  (1-x)(1 + x + x^2 + \cdots) = 1 + x + x^2 + \cdots - x - x^2 - \cdots = 1.
  \]
  Άρα, το $\sum_{n\ge0}x^n$ είναι αντιστρέψιμο με αντίστροφο $1-x$.

  Για το (γ) παρατηρούμε ότι $f(x)g(x) = 1$, για δύο τυπικές δυναμοσειρές $f(x) = \sum_{n\ge0}a_nx^n$ και $g(x) = \sum_{n\ge0}b_nx^n$ αν και μόνο αν 
  \begin{equation}
    \label{eq:FPS_1}
    a_0b_n + a_1b_{n-1} + \cdots + a_nb_0 = 
    \begin{cases}
    1, &\ \text{αν $n=0$} \\
    0, &\ \text{αν $n \ge 1$}.
    \end{cases}
  \end{equation}
  Επομένως, αν $f(x) \in \left(
      R[\![x]\!]
    \right)^\times$, τότε $a_0b_0 = 1$ και κατά συνέπεια $a_0 \in R^\times$. Αντιστρόφως, αν $a_0 \in R^\times$, τότε η εξίσωση \eqref{eq:FPS_1} έχει μοναδική λύση ως προς $(b_0, b_1, \dots)$, που δίνεται από 
    \begin{align*}
      b_0 &= a_0^{-1} \\
      b_1 &= -a_1b_0a_0^{-1} \\
      b_2 &= -(a_1b_1 + a_2b_0)a_0^{-1} \\
      &\ \ \vdots
    \end{align*}
\end{proof}

\begin{remark}
  Μπορούμε, λοιπόν, να γράψουμε 
  \[
  1 + x + x^2 + \cdots = \frac{1}{1-x},
  \]
  όπου το δεξί μέλος δεν είναι άλλο από το στοιχείο $(1-x)^{-1}$. 
\end{remark}

\begin{exercise}{\rm(\cite[Άσκηση~7.2.4]{DF04})}
  \label{ex:DF_7.2.4_7.4.18}
  Αν ο $R$ είναι ακέραια περιοχή, τότε να δείξετε ότι το ίδιο ισχύει και για τον $R[\![x]\!]$.
\end{exercise}

\begin{proof}[Λύση]
  Έστω $f(x)g(x) = 0$, για δύο τυπικές δυναμοσειρές $f(x) = \sum_{n\ge0}a_nx^n$ και $g(x) = \sum_{n\ge0}b_nx^n$ με $g(x) \neq 0$. Τότε 
  \begin{equation}
    \label{eq:FPS_2}
    a_0b_n + a_1b_{n-1} + \cdots + a_nb_0 = 0
  \end{equation}
  για κάθε $n \ge 0$. Για $n=0$, η Εξίσωση \eqref{eq:FPS_2} γίνεται $a_0b_0 = 0$. Επειδή ο $R$ είναι ακέραια περιοχή και $b_0 \ne 0$, προκύπτει ότι $a_0=0$. Για $n=1$, η Εξίσωση \eqref{eq:FPS_2} γίνεται $0 = a_0b_1 + a_1b_0 = a_1b_0$ κα γι αυτό $a_1 = 0$, όμοια με πριν. Συνεχίζοντας έτσι προκύπτει ότι $a_n = 0$, για κάθε $n \ge 0$. Με άλλα λόγια, ο δακτύλιος των τυπικών δυναμοσειρών δεν έχει μηδενοδιαιρέτες και γι αυτό είναι ακέραια περιοχή.
\end{proof}

% \begin{exercise}{\rm(\cite[Άσκηση~7.4.18]{DF04})}
%   \label{ex:DF_7.4.18}
%   Αν ο $R[\![x]\!]$ είναι ακέραια περιοχή, τότε να δείξετε ότι το ιδεώδες $(x)$ είναι πρώτο. Ειδικότερα, να δείξετε ότι το $(x)$ είναι μεγιστικό αν και μόνο αν το $R$ είναι σώμα.
% \end{exercise}

% \begin{proof}[Λύση]
%   Θεωρούμε την απεικόνιση 
%   \begin{align*}
%     \phi : R[\![x]\!] &\to R \\
%     \sum_{n\ge0}a_nx^n &\mapsto a_0.
%   \end{align*}
%   Αυτή είναι επιμορφισμός δακτύλιων με πυρήνα $\ker(\phi) = (x)$.   Από το πρώτο θεώρημα ισομορφισμού έπεται ότι $R \cong R[\![x]\!]/(x)$. Το ζητούμενο έπεται από τις Προτάσεις 1.11 και 1.10, αντίστοιχα.
% \end{proof}

Για την επόμενη άσκηση, θυμίζουμε ότι $F$-άλγεβρα είναι ένας δακτύλιος $A$ εφοδιασμένος με δομή διανυσματικού χώρου πάνω από το $F$ με τις δύο δομές να ικανοποιούν τη συνθήκη 
\[
\lambda(ab) = (\lambda{a})b = a(\lambda{b})
\]
για κάθε $a, b \in A$ και $\lambda \in F$. Για παράδειγμα, ο δακτύλιος των πολυωνύμων $F[x_1, x_2, \dots, x_n]$ είναι μια $F$-άλγεβρα (γιατί;).

\begin{exercise}
  \label{ex:hilbert_series}
Μια $F$-άλγεβρα $A$ ονομάζεται \defn{διαβαθμισμένη} (graded) αν υπάρχει διάσπαση 
\[
A = A_0 \oplus A_1 \oplus \cdots 
\]
διανυσματικών χώρων πάνω από το $F$, τέτοια ώστε $A_0 \cong F$ και 
\[
A_nA_m \subseteq A_{n+m}
\]
για κάθε $n, m \ge 0$. Αν $\dim_F(A_n) < \infty$, για κάθε $n \ge 0$, τότε η τυπική δυναμοσειρά 
\[
\Hilb(A;t) \coloneqq \sum_{n \ge 0} \dim_F(A_n) \, t^n \ \in \ \ZZ[\![t]\!]
\]
ονομάζεται \defn{σειρά Hilbert} (Hilbert series) της $A$. Να δείξετε ότι ο πολυωνυμικός δακτύλιος $F[x_1, x_2, \dots, x_n]$ είναι διαβαθμισμένη $F$-άλγεβρα και να υπολογίσετε την αντίστοιχη σειρά Hilbert.
\end{exercise}

Πριν την λύση της άσκησης, ας δούμε την περίπτωση όπου $n=1$. Έχουμε τη διάσπαση 
\[
F[x] = F\{1\} \oplus F\{x\} \oplus F\{x^2\} \oplus \cdots 
\]
ως διανυσματικοί χώροι πάνω από το $F$ και παρατηρούμε ότι 
\[
(ax^n)(bx^m) = abx^{n+m}
\]
για κάθε $a, b \in F$ και $n, m \ge 0$. Άρα, 
\[
F\{x^n\}F\{x^m\} = F\{x^{n+m}\},
\]
για κάθε $n, m \ge 0$ και γι αυτό ο $F[x]$ είναι διαβαθμισμένη $F$-άλγεβρα. Υπολογίζουμε 
\[
\Hilb(F[x]; t) = \sum_{n \ge 0} \dim_F(F\{x^n\}) \, t^n = \sum_{n\ge0}t^n = \frac{1}{1-t}.
\]

\begin{proof}[Λύση της Άσκησης \ref{ex:hilbert_series}]
  Ένα πολυώνυμο του $F[x_1,x_2, \dots, x_n]$ ονομάζεται \defn{ομογενές} αν κάθε μονώνυμο που εμφανίζεται στο ανάπτυγμά του έχει τον ίδιο βαθμό. Για παράδειγμα, το 
  \[
  x_1x_2x_3 + x_1^2x_2
  \] 
  είναι ομογενές πολυώνυμο του $F[x_1,x_2,x_3]$ βαθμού $3$, ενώ το 
  \[
  x_1^3 + x_1x_2
  \]
  όχι. Κάθε πολυώνυμο του $F[x_1, x_2, \dots, x_n]$ μπορεί να γραφεί (με μοναδικό) τρόπο ως άθροισμα ομογενών πολυωνύμων ομαδοποιώντας όλα τα μονώνυμα που εμφανίζονται στο ανάπτυγμά του ανάλογα με τον βαθμό τους (γιατί;).
  
  
  Aν $A_d$ είναι το σύνολο όλων των ομογενών πολυωνύμων βαθμού $d$, τότε 
  \[
  A_d = F\{x_1^{\alpha_1}x_2^{\alpha_2}\cdots x_n^{\alpha_n} : \alpha_1, \alpha_2, \dots, \alpha_n \in \NN \ \text{και} \ \alpha_1 + \alpha_2 + \cdots + \alpha_n = d\},
  \]
  Από την παραπάνω παρατήρηση έπεται ότι 
  \[
  F[x_1, x_2, \dots, x_n] = A_0 \oplus A_1 \oplus \cdots 
  \]
  ως διανυσματικός χώρος πάνω από το $F$. Επιπλέον, έχουμε $A_kA_\ell \subseteq A_{k+\ell}$, για κάθε $k, \ell \ge 0$. Πράγματι,  
  \[
  \left(a x_1^{\alpha_1}x_2^{\alpha_2}\cdots x_n^{\alpha_n}\right)
  \left(b x_1^{\beta_1}x_2^{\beta_2}\cdots x_n^{\beta_n}\right) = 
  abx_1^{\alpha_1+\beta_1}x_2^{\alpha_2+\beta_2}\cdots x_n^{\alpha_n+\beta_n} 
  \]
  για κάθε $a, b \in F$ και $(\alpha_1, \alpha_2, \dots, \alpha_n), (\beta_1, \beta_2, \dots, \beta_n) \in \NN^n$ με $\sum_{i=1}^n \alpha_i = k$ και $\sum_{i=1}^n \beta_i = \ell$, αντίστοιχα. Άρα, ο πολυωνυμικός δακτύλιος είναι διαβαθμισμένη $F$-άλγεβρα.

  Για τον υπολογισμό της σειράς Hilbert έχουμε 
  \begin{align*}
    \Hilb(F[x_1, x_2, \dots, x_n]; t) 
    &= \sum_{d \ge 0} \dim_F(A_d) \, t^d \\ 
    &= \sum_{d \ge 0} \abs{\{(\alpha_1, \alpha_2, \dots, \alpha_n) \in \NN^n : \alpha_1 + \alpha_2 + \cdots + \alpha_n = d\}} \, t^d \\
    &= \sum_{\alpha_1, \alpha_2, \dots, \alpha_n \in \NN} t^{\alpha_1 + \alpha_2 + \cdots + \alpha_n} \\ 
    &= \prod_{i=1}^n \left(
      \sum_{\alpha_i \in \NN} t^{\alpha_i} 
    \right)\\ 
    &= \prod_{i=1}^n \frac{1}{1-t} \\
    &= \frac{1}{(1-t)^n}.
  \end{align*}
\end{proof}

\begin{remark}
Μια ακολουθία $\alpha = (\alpha_1, \alpha_2, \dots, \alpha_n)$ μη αρνητικών ακεραίων τέτοια ώστε $\alpha_1 + \alpha_2 + \cdots + \alpha_n = d$ ονομάζεται \emph{ασθενής σύνθεση} (weak composition) του $d$ με $n$ μέρη. Για παράδειγμα, οι ασθενείς συνθέσεις του $2$ με $3$ μέρη είναι 
\[
  (2,0,0), \, (0,2,0), \, (0,0,2), \, (1,1,0), \, (1,0,1), \, (0,1,1)
\]
κτλ. Κάθε ασθενής σύνθεση του $d$ με $n$ μέρη αντιστοιχεί σε ένα μονώνυμο στις μεταβλητές $x_1, x_2, \dots, x_n$ βαθμού $d$. Στο παράδειγμα, τα αντίστοιχα μονώνυμα είναι 
\[
  x_1^2, \ x_2^2, \ x_3^2, \ x_1x_2, \ x_1x_3, \ x_2x_3.
\]
Επομένως, η διάσταση του $A_d$, δηλαδή το πλήθος των ομογενών πολυωνύμων βαθμού $d$ σε $n$ μεταβλητές ισούται με το πλήθος των ασθενών συνθέσεων του $d$ σε $n$ μέρη. Ισχυριζόμαστε ότι 
\[
\dim_F(A_d) = \binom{n + d - 1}{d}.
\]

Θα κατασκευάσουμε μια αμφιμονοσήμαντη αντιστοιχία μεταξύ του συνόλου των ασθενών συνθέσεων του $d$ με $n$ μέρη και του συνόλου των υποσυνόλων του $[n+d-1]$ με $n-1$ μέρη. Έστω $\alpha = (\alpha_1, \alpha_2, \dots, \alpha_n)$ μια ασθενής σύνθεση του $d$ με $n$ μέρη. Θέτουμε $\phi(\alpha) = \{s_1, s_2, \dots, s_{n-1}\}$, όπου 
\[
s_k = \alpha_1 + \alpha_2 + \cdots + \alpha_k + k.
\]
Παρατηρούμε ότι $1 \le s_1 < s_2 < \cdots < s_{n-1} = n+d-1$. Αντιστρόφως, αν $S = \{s_1, s_2, \dots, s_{n-1}\} \subseteq [n+d-1]$, τότε θέτουμε $\psi(T) = (\alpha_1, \alpha_2, \dots, \alpha_n)$, όπου 
\begin{align*}
  \alpha_1 &= s_1 - 1 \\
  \alpha_k &= s_k - s_{k-1} -1, \qquad \text{για $1 \le k \le n-1$} \\ 
  \alpha_n &= n + d - s_{n-1} - 1.
\end{align*}
Παρατηρούμε ότι $\alpha_k \ge 0$, για κάθε $1 \le k \le n$ και ότι 
\begin{align*}  
  \alpha_1 + \alpha_2 + \cdots + \alpha_n &= (s_1 - 1) + (s_2 - s_1 - 1) + \\ 
  &\quad \cdots + (s_{n-1} - s_{n-2} - 1) + (n + d - s_{n-1} - 1) \\
  &= n + d - n \\ 
  &= d.
\end{align*} 
Τέλος, οι απεικονίσεις $\phi$ και $\psi$ είναι αντίστροφες η μία της άλλης (γιατί;). 

Σε επίπεδο τυπικών δυναμοσειρών, η προηγούμενη άσκηση μας πληροφορεί ότι 
\[
\sum_{d \ge 0} \binom{n+d-1}{d} \, t^d = \frac{1}{(1-t)^n}
\]
στον $\ZZ[\![t]\!]$.
\end{remark}

% \begin{exercise}
%   Έστω $F$ ένα σώμα και $A$ μια διαβαθμισμένη $F$-άλγεβρα. Ένα ιδεώδες $I$ της $A$ ονομάζεται \defn{ομογενές} αν παράγεται από ομογενή πολυώνυμα. Να δείξετε ότι το $A/I$ είναι διαβαθμισμένη $F$-άλγεβρα με διάσπαση 
%   \[
%   A/I = A_0/I_0 \oplus A_1/I_1 \oplus \cdots.
%   \]
% \end{exercise}

% \begin{exercise}
%   Να υπολογιστεί η σειρά Hilbert του $R[x,y,z]/(xy)$.
% \end{exercise}

% \begin{proof}[Λύση]
  
% \end{proof}

\begin{thebibliography}{99}
    \bibitem{DF04}
    D. S. Dummit και R. M. Foote, 
    \textit{Abstract algebra},
    third edition, Wiley, 2004. 
    % \bibitem{Mal08}
    % Μ. Μαλιάκας, 
    % \textit{Εισαγωγή στην Μεταθετική Άλγεβρα}, 
    % Εκδόσεις Σοφία, 2008.
\end{thebibliography}
\end{document}